\chapter{Integrale complexe}



\section{Teorema lui Cauchy}
\index{integrală!complexă}
\indent\indent În multe situații, putem calcula integralele complexe direct,
într-o manieră asemănătoare cu integralele curbilinii. Un exemplu simplu:

\[
  I_1 = \int_{|z| = 1} z |dz|.
\]

Folosind forma polară, $ z = e^{it} $, deoarece integrala se face pe
$ |z| = 1 $, iar $ t \in [0, 2\pi] $. Rezultă $ dz = ie^{it} dt $,
deci $ |dz| = dt $. Atunci:
\[
  I_1 = \int_0^{2\pi} e^{it} dt = \dfrac{1}{i} e^{it} \Big|_0^{2\pi} = 0.
\]

Un alt exemplu:
\[
  I_2 = \int_S z |dz|,
\]
unde $ S $ este segmentul care unește pe $ 0 $ și $ i $. Putem parametriza
acest segment: $S : z = ti, t \in [0,1] $, deci $ dz = idt $ și din nou
$ |dz| = dt $. Rezultă:
\[
  I_2 = \int_0^1 tidt = \dfrac{1}{2}.
\]

Dar în unele situații, putem calcula chiar mai ușor:
\begin{theorem}[Cauchy]\label{thm:tcauchy} \index{teorema!Cauchy}
  Fie $ D \seq \CC $ un domeniu simplu conex și $ f : D \to \CC $ o funcție
  olomorfă pe $ D $, cu $ P = \dr{Re}f $ și $ Q = \dr{Im} f $, funcții de clasă
  $ \kal{C}^1(D) $.

  Fie $ \gamma : [a, b] \to D $ o curbă închisă și jordaniană (fără autointersecții)
  de clasă $ \kal{C}^1 $ pe porțiuni, astfel încît $ \dr{Int}\gamma $ să verifice
  condițiile formulei Green-Riemann.

  Atunci $ \ds\int_\gamma f(z) dz = 0 $.
\end{theorem}

Acesta este un caz simplu în care calculul se termină imediat cu rezultat nul.

În exerciții, vom folosi adesea următoarea:
\begin{theorem}[Formula integrală Cauchy]\label{thm:fintcauchy}\index{formula!Cauchy}
  Fie $ D \seq \CC $ un domeniu și $ f : D \to \CC $ o funcție olomorfă pe $ D $.
  Fie $ \overline{\Delta} \seq D $, unde $ \Delta $ este un domeniu simplu conex,
  mărginit, cu frontiera $ \gamma $, care este o curbă închisă, jordaniană, de
  clasă $ \kal{C}^1 $ pe porțiuni, orientată pozitiv.

  Atunci pentru orice $ a \in \Delta $ fixat are loc:
  \[
    f(a) = \dfrac{1}{2 \pi i} \int_\gamma \dfrac{f(z)}{z - a} dz.
  \]
\end{theorem}

Principala aplicație a acestei teoreme este să ne ajute să calculăm integrale pe
domenii în interiorul cărora funcția pe care o integrăm are probleme. Un exemplu:
\[
  \int_{|z - 2i| = 1} \dfrac{1}{z^2 + 4} dz.
\]
Observăm că $ z = 2i $ este un punct cu probleme pentru funcția considerată și
aplicăm formula integrală Cauchy.

Putem rescrie integrala astfel, izolînd punctul cu probleme:
\[
  \int_{|z - 2i| = 1} \dfrac{1}{z^2 + 4} dz = \int_{|z - 2i| = 1} \dfrac{\frac{1}{z + 2i}}{z - 2i} dz = %
  \dfrac{f(z)}{z - 2i} dz,
\]
unde am introdus exact funcția cu probleme, adică $ f(z) = \dfrac{1}{z + 2i} $.

Aplicăm formula integrală Cauchy și obținem:
\[
  f(2i) = \dfrac{1}{2\pi i} \int_{|z - 2i| = 1} \dfrac{f(z)}{z - 2i} dz \Rightarrow %
  \int_{|z - 2i| = 1} \dfrac{f(z)}{z - 2i} dz = 2\pi i f(2i) = \dfrac{\pi}{2}.
\]

Vor exista situații cînd punctul izolat nu poate fi eliminat atît de ușor (sau chiar deloc),
cazuri în care vom aplica un rezultat fundamental, \emph{teorema reziduurilor}.

%%%%%%%%%%%%%%%%%%%%%%%%%%%%%%%%%%%%%%%%%%%%%%%%%%%%%%%%%%%%%%%%%%%%%%

\section{Exerciții}

1. Calculați integrala $ \ds\int_{\Gamma} z^2 dz $, unde:
\begin{enumerate}[(a)]
\item $ \Gamma = [-1, i] \cup [i, 1] $;
\item $ \Gamma = \{ z(t) = 2 + it^2 \mid 0 \leq t \leq 1 \} $;
\item $ \Gamma = \{ z(t) = t + i \cos \dfrac{\pi t}{2} \mid -1 \leq t \leq 1 \} $;
\item $ \Gamma = OA $, cu $O(0, 0) $ și $ A(2, 1) $.
\end{enumerate}

\emph{Indicații:} Se parametrizează drumurile și se calculează ca în exemplele
de mai sus.

\vspace{1cm}

2. Folosind teorema Cauchy sau formula integrală Cauchy, calculați:
\begin{enumerate}[(a)]
\item $ \ds\int_{|z - 1| = 3} z^4 dz $;
\item $ \ds\int_{|z| = 4} \dfrac{\cos z}{z^2 - 6z + 5} $;
\item $ \ds\int_{|z| = 1} \dfrac{\sin z}{z(z-2)} dz $;
\item $ \ds\int_\Gamma \dfrac{\exp(z^2)}{z^2 - 6z} dz $, unde $ \Gamma: |z - 2| = r, r \in \{1, 3, 5\} $;
\item $ \ds\int_{|z| = 1} \dfrac{\exp(3z)}{z^4} dz $;
\item $ \ds\int_{|z - i| = 1} \dfrac{1}{(z^2 + 1)^2} dz $;
\item $ \ds\int_{|z| = 1} \dfrac{\sin z}{z^2(z - 2)} dz $.
\end{enumerate}

\emph{Indicații:} Ideea de bază este să identificăm punctele cu probleme ale
funcțiilor de integrat în interiorul domeniilor pe care integrăm, apoi să descompunem
integrandul cu o funcție căreia i se poate aplica teorema Cauchy.
\begin{enumerate}[(a)]
\item funcția $ z^4 $ este olomorfă, deci integrala este nulă;
\item avem $ \dfrac{\cos z}{(z - 1)(z - 5)} $, dar singurul punct cu probleme din interiorul
  domeniului este $ z_1 = 1 $. Definim $ f(z) = \dfrac{\cos z}{z - 5} $, iar integrala
  devine $ \ds\int_{|z| = 4} \dfrac{f(z)}{z - 1} dz $, care se calculează cu formula Cauchy.
\item pentru $ r = 1 $, funcția este olomorfă, deci integrala este nulă. Pentru $ r = 3 $,
  $ z = 0 $ este punct cu probleme, deci definim $ f(z) = \dfrac{\exp(z^2)}{z - 6} $.
\end{enumerate}

%%%%%%%%%%%%%%%%%%%%%%%%%%%%%%%%%%%%%%%%%%%%%%%%%%%%%%%%%%%%%%%%%%%%%%

\section{Teorema reziduurilor}

Similar cu orice funcție reală, și funcțiile complexe pot fi dezvoltate în serii
de puteri. În cazul complex, seriile se numesc \emph{serii Laurent} și pot conține
și puteri negative.

Informal, punctele cu probleme care ne interesează se numesc \emph{poli}
sau \emph{puncte singulare}. Ordinul unui pol $ z = a $ este multiplicitatea algebrică
a rădăcinii $ z = a $ în dezvoltarea în serie Laurent a funcției $ f $. În particular,
avem \emph{poli simpli, dubli} etc.

\begin{definition}\label{def:rez}\index{funcții complexe!reziduu}
  Fie $ f : \CC \to \CC $ o funcție complexă și dezvoltarea sa în serie Laurent
  în jurul unui punct $ z_0 \in \CC $:
  \[
    f(z) = \sum_{n \in \ZZ} a_n (z - z_0)^n, \quad a_n \in \RR.
  \]

  Se numește \emph{reziduul} funcției $ f $ în punctul singular $ z_0 $ coeficientul
  $ a_{-1} $ din dezvoltarea de mai sus, notat $ \mathrm{Rez}(f, z_0) $.
\end{definition}

Următoarea teoremă ne dă metode de calcul al reziduurilor, în funcție de multiplicitatea
lor:

\begin{theorem}[Calculul reziduurilor]\label{thm:calc-rez}\index{teorema!calculul reziduurilor}
  (1) $ \dr{Rez}(f, a) = c_{-1} $, unde $ c_{-1} $ este coeficientul lui $ \dfrac{1}{z-a} $
  în dezvoltarea în serie Laurent a funcției $ f $ în vecinătatea singularității $ z = a $.

  (2) Dacă $ z = a $ este pol de ordinul $ p \geq 2 $ pentru $ f $, atunci:
  \[
    \dr{Rez}(f, a) = \dfrac{1}{(p-1)!} \lim_{z \to a} \Big[ (z - a)^p f(z) \Big]^{(p-1)};
  \]

  (3) Dacă $ z = a $ este pol simplu pentru $ f $, atunci, particularizînd formula
  de mai sus, avem:
  \[
    \dr{Rez}(f, a) = \lim_{z \to a} (z - a) f(z);
  \]

  (4) Dacă $ f $ se poate scrie ca un cît de funcții, $ f = \dfrac{A}{B} $, olomorfe în
  jurul lui $ a $ și dacă $ z = a $ este pol simplu pentru $ f $, adică $ B(a) = 0 $,
  atunci:
  \[
    \dr{Rez}(f, a) = \lim_{z \to a} \dfrac{A(z)}{B'(z)}.
  \]
\end{theorem}

Rezultatul esențial al acestei secțiuni ne arată că, dacă integrăm o funcție
cu probleme, valoarea integralei este dată în mod esențial de reziduurile sale:

\begin{theorem}[Teorema reziduurilor]\label{thm:t-rez}\index{teorema!reziduurilor}
  Fie $ D \seq \CC $ un domeniu și $ f : D - \{ \alpha_1, \dots, \alpha_k \} \to \CC $ o funcție
  olomorfă pentru care $ \alpha_i $ sînt poli.

  Fie $ K \seq D $ un compact cu frontiera $ \Gamma = \partial K $, o curbă de clasă
  $ \kal{C}^1 $, jordaniană, orientată pozitiv și care conține toate $ \alpha_i $ în
  interior. Atunci:
  \[
    \int_\Gamma f(z) dz = 2 \pi i \sum_{j = 1}^k \dr{Rez}(f, \alpha_j).
  \]
\end{theorem}

De exemplu, să calculăm integrala:
\[
  I = \int_{|z| = r} \frac{e^z}{(z - i)(z - 2)} dz, r > 0, r \neq 1, 2.
\]

\emph{Soluție:} Dacă $ 0 < r < 1 $, putem aplica teorema lui Cauchy (\ref{thm:tcauchy})
și găsim $ I = 0 $.

Dacă $ 1 < r < 2 $, aplicăm formula integrală a lui Cauchy (\ref{thm:fintcauchy}) și găsim:
\[
  \int_{|z| = r} \frac{e^z}{(z - i)(z - 2)} dz = %
  \int_{|z| = r} \frac{\frac{e^z}{z - 2}}{z - i} dz = %
  \int_{|z| = r} \frac{f(z)}{z - i} dz,
\]
unde $ f(z) = \dfrac{e^z}{z - 2} $. Rezultă:
\[
  \int_{|z| = r} \frac{f(z)}{z - i} dz = 2 \pi i f(i) = 2 \pi i \frac{e^i}{i - 2}.
\]

Dacă $ r > 2 $, aplicăm teorema reziduurilor, cu $ i $ și $ 2 $ poli simpli. Avem:
\begin{align*}
  \mathrm{Rez}(f, i) &= \lim_{z \to i} (z - i) \frac{e^z}{(z - i)(z - 2)} = \frac{e^i}{i - 2} \\
  \mathrm{Rez}(f, 2) &= \lim_{z \to 2} \frac{e^z}{(z - i)(z - 2)} = \frac{e^2}{2 - i}.
\end{align*}

Rezultă, din teorema reziduurilor:
\[
  \int_{|z| = r} \frac{e^z}{(z - i)(z - 2)} dz = 2 \pi i \Big( \frac{e^i}{i - 2} + %
  \frac{e^2}{2 - i} \Big).
\]

%%%%%%%%%%%%%%%%%%%%%%%%%%%%%%%%%%%%%%%%%%%%%%%%%%%%%%%%%%%%%%%%%%%%%%

\section{Exerciții}

1. Calculați reziduurile funcțiilor în punctele $ a $ indicate:
\begin{enumerate}[(a)]
\item $ f(z) = \dfrac{\exp\big(z^2\big)}{z - 1}, a = 1 $;
\item $ f(z) = \dfrac{\exp\big(z^2\big)}{(z - 1)^2}, a = 1 $;
\item $ f(z) = \dfrac{z + 2}{z^2 - 2z}, a = 0 $;
\item $ f(z) = \dfrac{1 + e^z}{z^4}, a = 0 $;
\item $ f(z) = \dfrac{\sin z}{4z^2}, a = 0 $;
\item $ f(z) = \dfrac{z}{1 - \cos z}, a = 0 $.
\end{enumerate}

\emph{Indicație:} Verificăm multiplicitatea polului $ z = a $ și aplicăm formula
corespunzătoare din teorema \ref{thm:calc-rez}.

\vspace{1cm}

2. Să se calculeze următoarele integrale:
\begin{enumerate}[(a)]
\item $ I = \ds\int_{|z| = 2} \frac{dz}{z^2 - 1} $;
\item $ I = \ds\int_\gamma \frac{dz}{z^4 + 1}, \gamma: x^2 + y^2 - 2x = 0 $;
\item $ I = \ds\int_{|z| = 3} \frac{z^2 + 1}{(z - 1)^2 (z + 2)} dz $.
\end{enumerate}

\emph{Soluție:} (a) Punctele $ z = \pm 1 $ sînt poli de ordinul 1 pentru
funcția $ f(z) = \dfrac{1}{z^2 - 1} $. Ele sînt situate în interiorul discului pe
care integrăm, cu $ |z| = 2 $, deci putem aplica teorema reziduurilor:
\[
  I = 2 \pi i \cdot \Big( \mathrm{Rez}(f, z_1) + \mathrm{Rez}(f, z_2) \Big).
\]

Calculăm separat reziduurile:
\begin{align*}
  \Rez(f, z_1) &= \lim_{z \to 1} (z - 1) \cdot \frac{1}{z^2 - 1} = \frac{1}{2} \\
  \Rez(f, z_2) &= \lim_{z \to -1} (z + 1) \cdot \frac{1}{z^2 - 1} = -\frac{1}{2}.
\end{align*}

Rezultă:
\[
  I = \int_{|z| = 2} \frac{dz}{z^2 - 1} = 2 \pi i \cdot \Big( \frac{1}{2} - \frac{1}{2} \Big) = 0.
\]

\vspace{1cm}

(b) Curba $ \gamma $ este un cerc centrat în $ (1, 0) $ și cu raza $ 1 $. Căutăm
polii funcției $ f(z) = \dfrac{1}{z^4 + 1} $ care se află în interiorul lui $ \gamma $.

Avem succesiv:
\begin{align*}
  z^4 + 1 &= 0 \Rightarrow z^4 = -1 = \cos \pi + i \sin \pi \Rightarrow \\
  z &= \sqrt[4]{\cos \pi + i \sin \pi} \Rightarrow \\
  z_k &= \cos \frac{\pi + 2k\pi}{4} + i \sin \frac{\pi + 2k\pi}{4}, k = 0, 1, 2, 3.
\end{align*}

Doar punctele $ z_0, z_3 $ se află în interiorul discului delimitat de $ \gamma $
și calculăm reziduurile în aceste puncte.

Putem aplica formula din Teorema \ref{thm:calc-rez} (4) și avem:
\begin{align*}
  \Rez(f, z_0) &= \frac{A(z)}{B'(z)}\Big|_{z_0} = \frac{1}{4z^3}\Big|_{z_0} %
                 = -\frac{1}{4} e^{\frac{\pi i}{4}} \\
  \Rez(f, z_3) &= \frac{A(z)}{B'(z)}\Big|_{z_3} = \frac{1}{4z^3}\Big|_{z_3} %
                 = - \frac{1}{4} e^{\frac{7 \pi i}{4}}.
\end{align*}

Rezultă:
\[
  I = \int_\gamma \frac{dz}{z^4 + 1} = %
  2 \pi i \Big( -\frac{1}{4} e^{\frac{\pi i}{4}} - \frac{1}{4} e^{\frac{7 \pi i}{4}} \Big) = %
  -\frac{\pi \sqrt{2} i}{2}.
\]

\vspace{1cm}

(c) Avem doi poli, $ z = 1, z = -2 $ în interiorul conturului. Se vede că $ z_1 = 1 $
este pol de ordinul 2, iar $ z_2 = -2 $ este pol de ordinul 1. Calculăm reziduurile:
\begin{align*}
  \Rez(f, z_1) &= \frac{1}{2!} \lim_{z \to 1} \Big[ (z - 1)^2 \frac{z^2 + 1}{(z - 1)^2(z + 2)} \Big] \\
               &= \frac{1}{2!} \lim_{z \to 1} \frac{z^2 + 4z - 1}{(z + 2)^2} \\
               &= \frac{2}{9} \\
  \Rez(f, z_2) &= \lim_{z \to -2} (z + 2) \frac{z^2 + 1}{(z - 1)^2(z + 2)} = \frac{5}{9}.
\end{align*}

Rezultă:
\[
  I = \int_{|z| = 3} \frac{z^2 + 1}{(z - 1)^2(z + 2)} dz = \frac{14}{9} \pi i.
\]

\vspace{1cm}

3. Să se calculeze integralele:
\begin{enumerate}[(a)]
\item $ \ds\int_{|z| = 1} \frac{dz}{\sin z} $;
\item $ \ds\int_{|z| = 1} \frac{e^z}{z^2}dz $;
\item $ \ds\int_{|z| = 5} z e^{\frac{3}{z}} dz $;
\item $ \ds\int_{|z - 1| = 1} \frac{dz}{(z + 1)(z - 1)^3} $;
\item $ \ds\int_{|z| = 1} \frac{\sin z}{z^4} dz $.
\end{enumerate}

\emph{Indicații}: (a) $ z = 0 $ este singurul pol din interiorul domeniului;

(b), (c): Dezvoltăm în serie Laurent și identificăm reziduurile folosind definiția.

(d) Avem $ z = 1 $ pol de ordin 3 și $ z = -1 $ pol simplu. Doar $ z = 1 $ se află
în interiorul domeniului și dezvoltăm în serie Laurent după puterile lui $ z - 1 $:
\begin{align*}
  \frac{1}{z + 1} &= \frac{1}{2 - (-(z - 1))} \\
                  &= \frac{1}{2} \frac{1}{1 - \big(-\frac{z - 1}{2}\big)} \\
                  &= \frac{1}{2} \sum_{n \geq 0} (-1)^n \frac{(z - 1)^n}{2^n} \\
                  &= \sum_{n \geq 0} (-1)^n \frac{(z - 1)^n}{2^{n+1}},
\end{align*}
pentru $ |z - 1| < 2 $.

Rezultă:
\[
  \frac{1}{(z + 1)(z - 1)^3} = \sum_{n \geq 0} (-1)^n \frac{(z - 1)^{n-3}}{2^{n+1}} = %
  \frac{1}{2(z - 1)^3} - \frac{1}{4(z - 1)^2} + \frac{1}{8(z - 1)} - \frac{1}{16} + \dots,
\]
deci $ \Rez(f, 1) = \frac{1}{8} $.

%%%%%%%%%%%%%%%%%%%%%%%%%%%%%%%%%%%%%%%%%%%%%%%%%%%%%%%%%%%%%%%%%%%%%%

\section{Aplicații ale teoremei reziduurilor}

\indent\indent Putem folosi teorema reziduurilor pentru a calcula integrale trigonometrice de forma:
\[
  I = \int_0^{2 \pi} R(\cos \theta, \sin \theta) d \theta,
\]
unde $ R $ este o funcție rațională.

Facem schimbarea de variabilă $ z = e^{i\theta} $ și atunci, pentru $ \theta \in [0, 2\pi] $,
$ z $ descrie cercul $ |z| = 1 $, o dată, în sens direct.

Folosim formulele lui Euler:
\begin{align*}
  \cos \theta &= \frac{e^{i\theta} + e^{-i\theta}}{2} = \frac{1}{2} \Big( z + \frac{1}{z} \Big) \\
  \sin \theta &= \frac{e^{i\theta} - e^{-i\theta}}{2i} = \frac{1}{2i} \Big( z - \frac{1}{z} \Big).
\end{align*}

Atunci, dacă $ z = e^{i\theta} $, rezultă $ dz = ie^{i\theta} d\theta = izd\theta $,
iar integrala devine:\footnote{$\oint$ marchează o integrală pe un contur \emph{închis}}
\[
  I = \int_0^{2 \pi} R(\cos \theta, \sin \theta) d\theta = \oint_{|z| = 1} R_1(z) dz,
\]
unde:
\[
  R_1(z) = \frac{1}{iz} R\Big( \frac{z^2 + 1}{2z}, \frac{z^2 - 1}{2iz} \Big).
\]

Această funcție poate avea poli și deci putem folosi teorema reziduurilor. Dacă
$ a_1, \dots, a_n $ sînt polii din interiorul cercului unitate, avem:
\[
  I_1 = 2 \pi i \sum_{k \geq 1} \Rez(R_1, a_k).
\]

\vspace{1cm}

Să vedem cîteva exemple:
\begin{enumerate}[(a)]
\item $ \ds\int_0^{2 \pi} \frac{d\theta}{2 + \cos \theta} $;
\item $ \ds\int_0^{2 \pi} \frac{d\theta}{1 + 3 \cos^2 \theta} $;
\item $ \ds\int_0^{2 \pi} \frac{\sin \theta}{1 + i \sin \theta} d\theta $;
\item $ \ds\int_0^{2 \pi} \frac{d\theta}{1 + a \sin \theta}, |a| < 1, a \in \RR $.
\end{enumerate}

\emph{Soluție:}

(a) Notăm $ z = e^{i\theta} $, cu $ \theta \in [0, 2\pi] $. Atunci avem succesiv:
\begin{align*}
  dz &= ie^{i\theta} d\theta = izdz \Rightarrow d\theta = \frac{dz}{iz} \\
  \cos \theta &= \frac{e^{i\theta} + e^{-i\theta}}{2} = \frac{1}{2} \Big( z + \frac{1}{z} \Big) \\
  \frac{1}{2 + \cos \theta} &= \frac{2z}{z^2 + 4z + 1} \\
  \int_0^{2 \pi} \frac{d\theta}{2 + \cos \theta} &= \oint_{|z| = 1} \frac{2z}{z^2 + 4z + 1} \frac{dz}{iz} \\
     &= -2i \oint_{|z| = 1} \frac{dz}{z^2 + 4z + 1}.
\end{align*}

Acum folosim teorema reziduurilor. Singularitățile funcției $ f(z) = \dfrac{1}{z^2 + 4z + 1} $
sînt $ z = -2 \pm \sqrt{3} $, care sînt poli simpli. Numai $ z = -2 + \sqrt{3} $ se află
în interiorul cercului $ |z| = 1 $ și calculăm reziduul folosind
Teorema \ref{thm:calc-rez}(2).
\index{integrală!trigonometrică}

%%% Local Variables:
%%% mode: latex
%%% TeX-master: "../am2-18-19"
%%% End:
